\documentclass[11pt]{article}
\usepackage{icmlko}
\begin{document}

\title{판에 구멍이 있고, 그 구멍이 진짜 개선의 두 배다 --- 안 쓰고
있었을 뿐이다}
\author{월드모델 실험실}
\date{노트 446 · 2026년 8월 2일}
\maketitle

\begin{abstract}
노트 445 가 ``판의 여지는 시간에 있다''로 끝나서 시간을 모형에 넣어 보려고
코드를 열었다. 먼저 확인한 것 --- \textbf{챔피언은 \texttt{t} 를 아예 안
쓴다.} 모든 \texttt{predict} 가 \texttt{t} 를 받는데 \texttt{\_feat} 은
받지도 않는다. 2년 동안 출시일이 그대로 버려지고 있었다.

넣기 전에 유보 안을 봤다. \textbf{날짜와 라벨이 가중평균 $-0.228$ 로 붙어
있다}(도서 $-0.562$ · 웹툰 $-0.350$). 늦게 나온 것일수록 라벨이 낮다 ---
라벨이 관측 시점까지 \emph{쌓이는} 양이니 \textbf{늦게 나온 쪽이 기계적으로
낮은} 것이지 작품이 나빠서가 아니다. 판은 유보 위의 \emph{날} rho 라
날짜를 통제하지 않는다. \textbf{구멍이다.}

그래서 넣는 대신 \textbf{구멍이 얼마짜리인지}를 쟀다. 최근순 축 하나를
넣고, 판을 두 가지로 찍었다.

\begin{center}
\begin{tabular}{lrr}
\hline
 & 판 (날것) & 날짜 통제 판 \\
\hline
챔피언 & $+0.4764$ & $+0.5158$ \\
최근순 축 하나 추가 & $\mathbf{+0.4950}$ & $\mathbf{+0.4864}$ \\
\hline
차 & $\mathbf{+0.0186}$ & $\mathbf{-0.0294}$ \\
\hline
\end{tabular}
\end{center}

\textbf{날것은 오르고 실력은 내린다.} 그리고 크기를 보라 --- 이 사이클에서
실제로 채택된 것 중 제일 큰 것이 깊이 $4{\to}6$ 의 $+0.0093$ 이다.
\textbf{구멍이 진짜 개선의 두 배다.}
\end{abstract}

\begin{figure}[t]
\centering
\includegraphics[width=\textwidth]{figs/f1.pdf}
\caption{\textbf{가} --- 축 하나로 날 판은 오르고 날짜 통제 판은 내린다.
\textbf{나} --- 도메인이 받는 양은 그 도메인에서 날짜가 라벨에 얼마나
붙어 있는지로 갈린다. 도서($-0.562$)가 혼자 $+0.2376$ 을 가져간다.}
\end{figure}

\section{챔피언은 왜 안 걸렸나}

운이었다. 챔피언의 예보와 날짜의 상관은 $-0.067$ 이고, 날짜를 통제해도
판이 $+0.5143 \to +0.5158$ 로 \emph{안 내려간다} --- 지금 판 점수 중
날짜에서 온 것이 사실상 없다. \textbf{구멍이 없어서가 아니라 아무도 그
축을 안 넣어서다.}

노트 445 의 결론(``여지는 시간에 있다'')을 곧이곧대로 따랐으면 이번
사이클에 최근순 축을 넣고 $+0.0186$ 을 채택했을 것이다. 판은 통과하고
모형은 나빠진다.

\section{규약 --- 판이 오르는데 날짜 통제 판이 내려가면 안 넣는다}

순위 규약은 \textbf{안 바꾼다}. 순위는 그대로 유보 표본의 판 rho 다 ---
지표를 바꾸는 것은 이 노트가 할 일이 아니고, 바꾸면 지난 2년 점수와
못 견준다. 대신 \textbf{거부권 있는 가드}를 하나 붙인다.

\begin{quote}
\textbf{판이 오르는데 날짜 통제 판이 내려가는 변경은 판이 통과해도 안
넣는다.}
\end{quote}

노트 420 의 ``전이를 확실히 잃는 변경은 판이 통과해도 안 넣는다''와 같은
꼴이다. 자료 쪽 절반은 \texttt{lab/loop.py} 의 \texttt{\_date\_hole} 로
모든 실행에 자동으로 찍는다 --- 도메인별 유보 안 날짜↔라벨 상관과 구멍이
큰 도메인 목록(지금은 게임 · 도서 · 만화 · 모바일 · 웹툰 다섯), 그리고
\textbf{빠진 도메인이 왜 빠졌는지}(노트 401 규약: 시장팝업 · 아이돌 ·
펀딩은 유보 날짜가 한두 값뿐이라 못 잰다).

\section{이것이 노트 445 를 어떻게 고치나}

노트 445 는 옳았다 --- 예보 값 $-0.1492$ 가 모음 값 $+0.0699$ 보다 크다.
\textbf{다만 그 여지는 판 위에서는 정직하게 못 딴다.} 앞을 내다보는 일이
어려운 것과, 앞을 내다보는 척하며 \emph{언제 나왔는지}를 읽는 것은 다르다.
시간을 모형의 재료로 쓰려면 \textbf{날짜 통제 판이 같이 올라야} 한다 ---
계절성 · 출시 간격 · 경쟁 밀도 같은 것이지 \emph{최근순}이 아니다.

\section*{남은 것}

\begin{itemize}
\item \textbf{정직한 시간 축.} 출시 요일 · 월 · 그 시기 같은 도메인의
출시 밀도 --- 최근순과 달리 유보 안에서 라벨과 기계적으로 안 붙는 것들.
날짜 통제 판으로 잰다.
\item \textbf{라벨의 쌓임을 직접 고칠 수 있나.} 라벨을 출시 후 고정
기간(예: 90일)으로 자르면 구멍이 원천에서 닫힌다. 노트 210 이 게임에서
이미 30일 창을 했다 --- 다른 도메인도 되나. 이건 수집 갈래다.
\item 지난 2년 점수 중 \textbf{날짜에서 온 것이 있나} --- 챔피언은
$-0.067$ 로 깨끗하지만 중간에 죽은 도전자들은 안 찍어 봤다.
\end{itemize}

\end{document}
